\documentclass{article}
\usepackage{geometry}
\usepackage{graphicx} % Required for inserting images
\usepackage{amsmath, amsthm, amssymb}
\usepackage{parskip}
\newgeometry{vmargin={15mm}, hmargin={24mm,34mm}}
\theoremstyle{definition} 
\newtheorem{definition}{Definition}

\newtheorem{theorem}{Theorem}[section]
\newtheorem{lemma}[theorem]{Lemma}
\newtheorem{proposition}[theorem]{Proposition}
\newtheorem{example}[theorem]{Example}
\newtheorem{corollary}{Corollary}[theorem]
\newcommand{\nil}{\operatorname{nil}}
\newcommand{\ann}{\operatorname{Ann}}

\title{ATOCA Exercises}
\date{April 2025}
\author{Boran Erol}

\begin{document}

\maketitle

\begin{lemma}[ATOCA 3.38]
    Let $ R $ be a ring and $ I,J $ be ideals of $ R $.
    Assume $ J $ is finitely generated modulo $ I $ and
    assume that $ J \subseteq \sqrt{I} $. Then, there exists some
    $ n \geq 1 $ such that $ J^{n} \subseteq I $.
\end{lemma}
\begin{proof}
    If $ J $ is contained $ I $, this is trivial. So assume $ J $ is not contained in $ I $.

    Let $ \bar{x}_{1}, \ldots, \bar{x}_{n} $ be the generators of $ J $ in $ R/I $.

    Since for all $ i \in [n] $ we have that $ x_{i} \in J $, for all $ i \in [n] $
    there exists some $ k_{i} \geq 1 $ such that $ x_{i}^{k_{i}} \in I $.

    Let $ k = k_{1} + \cdots + k_{n} $. By the binomial formula, $ J^{k} $ is contained
    in $ \langle x_{1}^{k}, x_{2}^{k}, \ldots, x_{n}^{k} \rangle $.

    But $ \langle x_{1}^{k}, x_{2}^{k}, \ldots, x_{n}^{k} \rangle $ is the zero ideal in $ R/I $
    by construction, so $ J^{k} $ is also the zero ideal, and we have the desired result.
\end{proof}

\end{document}
